% ============================================================
% UNIDAD 11: ENSAYOS, CALIBRACIÓN Y DIAGNÓSTICO DE FALLAS
% ============================================================
\clearpage
\section*{UNIDAD 11: ENSAYOS Y DIAGNÓSTICO}
\subsection*{Sistema pitot-estático, calibración, chequeo funcional y fallas típicas}

\begin{recuadro}
\textbf{Objetivo de la unidad:} realizar los ensayos funcionales de los instrumentos integrados
instalados en la aeronave, en particular el ensayo del sistema pitot-estático, y diagnosticar las
fallas más frecuentes, registrando el trabajo según los procedimientos aeronáuticos.
\end{recuadro}

% ------------------------------------------------------------
\section*{1. ENSAYO DEL SISTEMA PITOT-ESTÁTICO}

Se realiza con un \textbf{banco de prueba pitot-estático} (manual o automático), que se conecta a
las tomas pitot y estática con adaptadores y permite simular altitudes y velocidades aplicando
vacío y presión de forma controlada.

\subsubsection*{1.1. Precauciones}
\begin{itemize}
  \item Leer el procedimiento del AMM y del fabricante del equipo de prueba.
  \item \textbf{Nunca aplicar presión en el sistema estático ni vacío en el pitot} (salvo lo que
        indique el procedimiento) y respetar las velocidades máximas de cambio: los cambios bruscos
        dañan los sensores y los instrumentos.
  \item Mantener siempre la presión del pitot \textbf{mayor o igual} que la estática.
  \item Tapar las tomas no usadas con los adaptadores del banco y retirarlos todos al terminar (con
        cintas de aviso ``REMOVER ANTES DEL VUELO'').
\end{itemize}

\subsubsection*{1.2. Ensayo de pérdidas}
\begin{itemize}
  \item \textbf{Sistema estático:} en aeronaves no presurizadas, la norma de la FAA (FAR 43,
        apéndice E) indica evacuar el sistema hasta que el altímetro marque unos \textbf{1000 pies
        por encima} de la elevación del lugar; sin seguir bombeando, en \textbf{un minuto} la
        altitud indicada \textbf{no debe bajar más de 100 pies}.
  \item \textbf{Sistema pitot:} se aplica presión hasta una velocidad indicada de prueba y se
        verifica que la pérdida en un tiempo dado no supere lo que indica el manual de la
        aeronave.
  \item Las pérdidas típicas aparecen en mangueras resecas, uniones, drenajes y en los equipos
        recién instalados (conexiones de pitot y estática detrás del panel).
\end{itemize}

\subsubsection*{1.3. Verificación de altímetros, velocímetros y ADC}
Con el banco se llevan los instrumentos a distintos valores (altitudes y velocidades de prueba) y
se comparan con las tolerancias del manual. En los sistemas digitales se verifica también la
\textbf{altitud que el ADC envía al transpondedor} (correlación) y que todas las pantallas indiquen
lo mismo. Las normas de la FAA (FAR 91.411 y 91.413) exigen estos ensayos cada 24 meses calendario
para vuelo IFR en espacio controlado y para el transpondedor; verificar los requisitos de las RAAC
vigentes.

% ------------------------------------------------------------
\section*{2. OTROS CHEQUEOS FUNCIONALES}
\begin{center}
\small
\begin{tabular}{|p{0.26\linewidth}|p{0.62\linewidth}|}
\hline
\textbf{Sistema} & \textbf{Chequeo} \\ \hline
AHRS & Alineación al encender, indicación de actitud con la aeronave nivelada, respuesta a
movimientos, ausencia de banderas. \\ \hline
Rumbo / magnetómetro & Comparación con rumbos conocidos (rosa de los vientos) y con la brújula;
verificación con cargas eléctricas encendidas y apagadas. \\ \hline
Indicación de motor & Simular sensores (década de resistencias, calibrador de termocuplas,
generador de pulsos) y verificar lecturas, colores y alertas. \\ \hline
Combustible & Verificar la indicación en varios niveles de carga y la alerta de combustible
bajo. \\ \hline
Interfaces & Navegación del GPS y del VOR/ILS en el HSI, comandos al piloto automático, altitud al
transpondedor, datos de tránsito y meteorología. \\ \hline
Alimentación & Funcionamiento con la batería de reserva al cortar la alimentación principal; prueba
de capacidad de la batería. \\ \hline
Iluminación y pantallas & Brillo, píxeles defectuosos, ventiladores. \\ \hline
\end{tabular}
\end{center}

% ------------------------------------------------------------
\section*{3. DIAGNÓSTICO DE FALLAS}

\paragraph*{Método.} Leer la falla reportada; reproducirla; revisar los registros y las páginas de
estado del sistema; seguir el procedimiento de aislamiento del fabricante; verificar
\textbf{alimentación, masas, conectores y cableado} antes de cambiar una unidad; reemplazar o
reparar; ensayar; registrar.

\begin{center}
\small
\begin{tabular}{|p{0.36\linewidth}|p{0.52\linewidth}|}
\hline
\textbf{Síntoma} & \textbf{Causas posibles} \\ \hline
Cruz roja en actitud y rumbo & AHRS sin alimentación o sin comunicación; conector; falla del AHRS
\\ \hline
Cruz roja en velocidad, altitud y variómetro & ADC; obstrucción o agua en pitot/estática (con
datos erróneos más que cruz roja); comunicación \\ \hline
Rumbo erróneo que cambia al encender equipos & Magnetómetro afectado por un campo magnético
cercano; calibración \\ \hline
Velocidad o altitud erróneas & Pérdidas, agua u obstrucciones en las líneas; tomas estáticas
dañadas \\ \hline
Indicación de motor errática & Masa del sensor, conector, cable de termocupla dañado o empalmado
con cable común \\ \hline
Pantalla que se reinicia & Caída de tensión en el arranque, mala conexión de alimentación,
sobretemperatura \\ \hline
Mensajes de incompatibilidad de configuración & Unidad reemplazada sin cargar software o
configuración \\ \hline
Sin navegación en el HSI & Configuración de la interfaz (ARINC 429, velocidad), cables A/B
invertidos \\ \hline
\end{tabular}
\end{center}

% ------------------------------------------------------------
\section*{4. REGISTRO DEL TRABAJO}
{\small
\renewcommand{\arraystretch}{1.45}
\begin{tabular}{|p{0.955\linewidth}|}
\hline
\textbf{TARJETA DE TRABAJO -- INSTRUMENTOS INTEGRADOS} \hfill N°: \makebox[2cm]{\dotfill} \\ \hline
Aeronave (matrícula y modelo): \dotfill \quad Fecha: \makebox[3cm]{\dotfill} \\ \hline
Equipo / sistema y referencia del manual: \dotfill \\ \hline
Falla reportada o tarea: \\[0.8cm] \hline
Acciones realizadas: \\[1.1cm] \hline
Unidad removida P/N y S/N: \dotfill \quad Instalada P/N y S/N: \makebox[3.6cm]{\dotfill} \\ \hline
Software y configuración cargados (versión): \dotfill \\ \hline
Ensayos realizados (pitot-estático, funcionales, calibraciones) y resultados: \\[1cm] \hline
Equipos de prueba y vencimiento de calibración: \dotfill \\ \hline
Ejecutó (firma y licencia): \dotfill \quad Inspeccionó: \makebox[4.3cm]{\dotfill} \\ \hline
\end{tabular}
}

% ------------------------------------------------------------
\section*{5. ACTIVIDADES DE LA UNIDAD}
\begin{enumerate}
  \item Realizar con el banco del taller un ensayo de pérdidas del sistema estático de una
        aeronave o de un tablero didáctico, y registrar el resultado.
  \item Verificar un altímetro a varias altitudes de prueba y comparar con las tolerancias.
  \item Simular con una década de resistencias y un calibrador las lecturas de temperatura de
        aceite y de un cilindro, y verificar la indicación.
  \item Resolver tres casos de falla planteados por el docente con el método de diagnóstico y
        completar la tarjeta de trabajo.
\end{enumerate}
